All graphs are finite, simple, undirected.
A \emph{diagonal} of a cycle $C=v_0v_1\dots v_{\ell-1}v_0$ means a chord: an edge $v_iv_j$ with $|i-j|\not\equiv 1 \pmod \ell$.

Question (1): If $G$ is $K_4$-free and $\chi(G)=4$, must $G$ contain an \emph{odd} cycle with at least two diagonals?

Question (2): Does there exist a function $f(r)\to\infty$ such that every graph $G$ with $\chi(G)=4$ and with the property that every subgraph on $\le r$ vertices has chromatic number $\le 3$ contains an odd cycle with at least $f(r)$ diagonals?
